\chapter{代数结构}
\intro{代数结构是数学抽象化的核心——将具体运算（加法、乘法、函数复合等）抽象为集合上的二元运算，研究其公理化性质。}

\section{代数系统的基本概念}

\subsection{二元运算}

\begin{mydefinition}{}
设 $S$ 是一个非空集合。一个\textbf{二元运算} $*$ 是从 $S \times S$ 到 $S$ 的一个映射：
\[
* : S \times S \to S, \quad (a,b) \mapsto a * b
\]
\end{mydefinition}

\textbf{运算的基本性质}：

\begin{center}
\begin{tabular}{c|c}
\hline
\textbf{性质} & \textbf{定义} \\
\hline
封闭性 & $\forall a,b\in S,\; a*b \in S$（二元运算定义内含封闭性）\\
结合律 & $\forall a,b,c\in S,\; (a*b)*c = a*(b*c)$ \\
交换律 & $\forall a,b\in S,\; a*b = b*a$ \\
\hline
\end{tabular}
\end{center}

\begin{example}
$(\mathbb{N}, +)$：封闭、结合、交换。$(\mathbb{N}, -)$：不封闭。$M_n(\mathbb{R})$ 上矩阵乘法：封闭、结合，但\textbf{不交换}。
\end{example}

\subsection{特殊元素：幺元、零元与逆元}

\begin{mydefinition}{幺元（单位元）}
元素 $e \in S$ 称为\textbf{幺元}，若 $\forall a \in S,\; e*a = a*e = a$。
如果幺元存在，则幺元唯一。
\end{mydefinition}

\begin{mydefinition}{零元}
元素 $\theta \in S$ 称为\textbf{零元}，若 $\forall a \in S,\; \theta*a = a*\theta = \theta$。
如果零元存在，则零元唯一。
\end{mydefinition}

\begin{mydefinition}{逆元}
设 $(S,*)$ 有幺元 $e$。对于 $a \in S$，若存在 $b \in S$ 使得 $a*b = b*a = e$，则 $b$ 称为 $a$ 的\textbf{逆元}，记作 $a^{-1}$。若运算满足结合律且有幺元，则每个元素最多有一个逆元。
\end{mydefinition}

\begin{proof}[逆元的唯一性]
设 $b, c$ 都是 $a$ 的逆元：
\[
b = b*e = b*(a*c) = (b*a)*c = e*c = c
\]
\end{proof}

\subsection{代数结构层次}

\begin{center}
\begin{tabular}{c|c|c}
\hline
\textbf{结构} & \textbf{条件} & \textbf{英文} \\
\hline
广群 (Magma)    & 封闭性 & Magma / Groupoid \\
半群 (Semigroup) & 封闭 + 结合律 & Semigroup \\
幺半群 (Monoid)   & 半群 + 幺元 & Monoid \\
群 (Group)       & 幺半群 + 每个元素有逆元 & Group \\
交换群 (Abel 群)  & 群 + 交换律 & Abelian Group \\
\hline
\end{tabular}
\end{center}

\begin{center}
\begin{tikzpicture}[
    dmAlgStruct/.style={rectangle, draw=blue!60, fill=blue!5, rounded corners, minimum width=2.8cm, align=center, minimum height=0.8cm},
    dmAlgInclude/.style={->, >=Stealth, thick, draw=blue!50}
]
\node[dmAlgStruct, fill=red!10, draw=red!60] (magma) at (0,0)        {广群 (Magma)\\[1pt] \tiny 封闭性};
\node[dmAlgStruct, above=0.6cm of magma] (semi)    {半群 (Semigroup)\\[1pt] \tiny 封闭 + 结合};
\node[dmAlgStruct, above=0.6cm of semi] (monoid)   {幺半群 (Monoid)\\[1pt] \tiny 半群 + 幺元};
\node[dmAlgStruct, above=0.6cm of monoid] (group)   {群 (Group)\\[1pt] \tiny 幺半群 + 逆元};
\node[dmAlgStruct, above=0.6cm of group] (abel)     {Abel 群\\[1pt] \tiny 群 + 交换律};

\draw[dmAlgInclude] (magma) -- (semi);
\draw[dmAlgInclude] (semi) -- (monoid);
\draw[dmAlgInclude] (monoid) -- (group);
\draw[dmAlgInclude] (group) -- (abel);

\node at (4, 2.8) [text width=3cm, align=left, font=\small\fangsong] {%
    \begin{tabular}{l}
    环 $\supset$ 整环 $\supset$ 域 \\
    环 $\supset$ 除环 $\supset$ 域
    \end{tabular}
};
\end{tikzpicture}
\end{center}

\section{半群与幺半群}

\subsection{半群}

\begin{mydefinition}{}
$(S,*)$ 是\textbf{半群}当且仅当 $*$ 封闭且满足结合律。
\end{mydefinition}

常见例子：$(\mathbb{N}, +)$、$(\mathbb{N}, \times)$、$(M_n(\mathbb{R}), \times)$（矩阵乘法）、$(\Sigma^+, \circ)$（字符串拼接）、$(\{0,1\}, \land)$（布尔运算）。

\begin{mydefinition}{循环半群}
半群中，若存在元素 $a$ 使得 $S = \{a^n \mid n \in \mathbb{Z}^+\}$，则称 $(S,*)$ 为\textbf{循环半群}，$a$ 称为生成元。循环半群总是交换的。
\end{mydefinition}

\subsection{幺半群}

\begin{mydefinition}{}
$(M,*)$ 是\textbf{幺半群}当且仅当 $(M,*)$ 是半群且存在幺元 $e \in M$。
\end{mydefinition}

常见幺半群：
\[
\begin{array}{c|c|c}
\hline
\text{幺半群} & \text{运算} & \text{幺元} \\
\hline
(\mathbb{N}_0, +) & 加法 & 0 \\
(\mathbb{N}, \times) & 乘法 & 1 \\
(\Sigma^*, \circ) & 字符串拼接 & 空串\ \varepsilon \\
(\mathcal{P}(X), \cup) & 集合并 & \emptyset \\
(M_n(\mathbb{R}), \times) & 矩阵乘法 & 单位阵 I_n \\
\hline
\end{array}
\]

\begin{mydefinition}{可逆元}
幺半群 $M$ 中，元素 $a$ 若存在 $b$ 使得 $a*b = b*a = e$，则 $a$ 称为\textbf{可逆元}。幺半群中所有可逆元的集合构成一个群。
\end{mydefinition}

\section{群}

\subsection{群的定义}

\begin{mydefinition}{}
代数系统 $(G,*)$ 称为\textbf{群}，当满足：
\begin{enumerate}
    \item \textbf{封闭性}：$\forall a,b \in G,\; a*b \in G$
    \item \textbf{结合律}：$\forall a,b,c \in G,\; (a*b)*c = a*(b*c)$
    \item \textbf{幺元存在}：$\exists e \in G,\; \forall a \in G,\; e*a = a*e = a$
    \item \textbf{逆元存在}：$\forall a \in G,\; \exists a^{-1} \in G,\; a*a^{-1} = a^{-1}*a = e$
\end{enumerate}
若还满足交换律 $a*b = b*a$，则称为\textbf{交换群}（Abel 群）。
\end{mydefinition}

\subsection{群的基本性质}

\begin{myTheorem}{消去律}
在群 $(G,*)$ 中，对任意 $a,b,c \in G$：
\begin{enumerate}
    \item 左消去律：若 $a*b = a*c$，则 $b = c$。
    \item 右消去律：若 $b*a = c*a$，则 $b = c$。
\end{enumerate}
\end{myTheorem}

\begin{myTheorem}{方程可解性}
在群 $(G,*)$ 中，对任意 $a,b \in G$：
\begin{itemize}
    \item 方程 $a*x = b$ 有唯一解 $x = a^{-1}*b$。
    \item 方程 $y*a = b$ 有唯一解 $y = b*a^{-1}$。
\end{itemize}
\end{myTheorem}

\begin{myTheorem}{}
群中恒有：$(a^{-1})^{-1} = a$，且 $(a*b)^{-1} = b^{-1}*a^{-1}$。幺元是唯一的幂等元。
\end{myTheorem}

\subsection{子群与判定定理}

\begin{mydefinition}{}
群 $(G,*)$ 的非空子集 $H$，若 $(H,*)$ 也构成群，则 $H$ 称为 $G$ 的\textbf{子群}，记作 $H \leq G$。$\{e\}$ 和 $G$ 本身称为\textbf{平凡子群}。
\end{mydefinition}

\begin{myTheorem}{子群判定定理 I：一步检验}
$G$ 的非空子集 $H$ 是子群当且仅当：
\[
\forall a,b \in H,\; a*b^{-1} \in H
\]
\end{myTheorem}

\begin{myTheorem}{子群判定定理 II：两步检验}
$H \leq G$ 当且仅当 (1) $H$ 非空;(2) $\forall a,b\in H, a*b\in H$;(3) $\forall a\in H, a^{-1}\in H$。
\end{myTheorem}

\begin{myTheorem}{有限子群判定}
若 $H$ 是 $G$ 的有限非空子集且对运算封闭，则 $H \leq G$。
\end{myTheorem}

\subsection{循环群}

\begin{mydefinition}{}
群 $G$ 若存在元素 $a \in G$ 使得 $G = \{a^k \mid k \in \mathbb{Z}\}$，则称 $G$ 为\textbf{循环群}，$a$ 称为生成元，记作 $G = \langle a \rangle$。
\end{mydefinition}

\begin{mydefinition}{元素的阶}
群中元素 $a$ 的\textbf{阶} $|a|$ 是满足 $a^k = e$ 的最小正整数 $k$。若不存在这样的 $k$，则称 $a$ 的阶为无穷大。
\end{mydefinition}

\textbf{性质}：
\begin{enumerate}
    \item 若 $|a| = n$，则 $\langle a \rangle = \{e, a, a^2, \dots, a^{n-1}\}$，$|\langle a \rangle| = n$。
    \item $a^k = e$ 当且仅当 $n \mid k$（$n$ 为 $a$ 的阶）。
    \item $a^k$ 的阶为 $\frac{n}{\gcd(n,k)}$;$a^k$ 是生成元当且仅当 $\gcd(n,k) = 1$。
    \item 循环群的子群仍然是循环群。
    \item 对 $n$ 阶循环群，对 $n$ 的每个正因子 $d$，恰有一个 $d$ 阶子群。
\end{enumerate}

重要例子：$(\mathbb{Z}, +)$（无限循环群，生成元为 $1$ 和 $-1$）;$(\mathbb{Z}_n, \oplus_n)$（$n$ 阶循环群，生成元有 $\varphi(n)$ 个）;$n$ 次单位根群（复数乘法下是循环群）。

\begin{center}
\begin{tikzpicture}[
    dmCayleyNode/.style={circle, draw=green!50!black, fill=green!10, minimum size=8mm, inner sep=0pt},
    dmCEdge/.style={->, >=Stealth, thick, draw=blue!60}
]
% Cayley diagram of Z_6 with generator 1
\node[dmCayleyNode] (c0) at (0:1.5cm) {0};
\node[dmCayleyNode] (c1) at (60:1.5cm) {1};
\node[dmCayleyNode] (c2) at (120:1.5cm) {2};
\node[dmCayleyNode] (c3) at (180:1.5cm) {3};
\node[dmCayleyNode] (c4) at (240:1.5cm) {4};
\node[dmCayleyNode] (c5) at (300:1.5cm) {5};
\draw[dmCEdge] (c0) -- (c1);
\draw[dmCEdge] (c1) -- (c2);
\draw[dmCEdge] (c2) -- (c3);
\draw[dmCEdge] (c3) -- (c4);
\draw[dmCEdge] (c4) -- (c5);
\draw[dmCEdge] (c5) -- (c0);
\node at (0, -2.2) {$\mathbb{Z}_6$ 的 Cayley 图（生成元 1）};
\end{tikzpicture}
\end{center}

\subsection{陪集与拉格朗日定理}

\begin{mydefinition}{左陪集}
设 $H \leq G$，$a \in G$。则 $aH = \{a*h \mid h \in H\}$ 称为 $H$ 在 $G$ 中的\textbf{左陪集}。
\end{mydefinition}

\textbf{性质}：
\begin{enumerate}
    \item $|aH| = |H|$（双射 $h \mapsto ah$）。
    \item 任意两个左陪集要么相等，要么不相交。
    \item 所有左陪集的并等于 $G$。因此左陪集构成 $G$ 的一个\textbf{划分}。
\end{enumerate}

\begin{myTheorem}{拉格朗日定理}
若 $G$ 是有限群，$H \leq G$，则
\[
\boxed{|H| \;\big|\; |G|}
\]
且 $[G:H] = |G|/|H|$ 称为 $H$ 在 $G$ 中的\textbf{指数}。
\end{myTheorem}

\begin{proof}
左陪集构成 $G$ 的划分，每个陪集大小等于 $|H|$。设共有 $k$ 个不同左陪集，则 $|G| = k \cdot |H|$，即 $|H|$ 整除 $|G|$。
\end{proof}

\begin{corollary}
有限群中任意元素 $a$ 的阶 $|a|$ 整除 $|G|$。
\end{corollary}

\begin{corollary}
素数阶群必为循环群（从而必为 Abel 群）。
\end{corollary}

\begin{center}
\begin{tikzpicture}[
    dmSGLatticeNode/.style={rectangle, draw=blue!60, fill=blue!5, rounded corners, minimum width=2cm, align=center, minimum height=0.7cm},
    dmSGLatLine/.style={thick, draw=blue!50}
]
% Subgroup lattice of Z_12 as an example
\node[dmSGLatticeNode, fill=red!10, draw=red!60] (g) at (0,5) {$\mathbb{Z}_{12}$（12 阶）};
\node[dmSGLatticeNode] (h4) at (-2.5,3.2) {$\langle 3 \rangle \cong \mathbb{Z}_4$};
\node[dmSGLatticeNode] (h6) at (0,3.2) {$\langle 2 \rangle \cong \mathbb{Z}_6$};
\node[dmSGLatticeNode] (h2) at (-2.5,1.4) {$\langle 6 \rangle \cong \mathbb{Z}_2$};
\node[dmSGLatticeNode] (h3) at (0,1.4) {$\langle 4 \rangle \cong \mathbb{Z}_3$};
\node[dmSGLatticeNode, fill=green!10, draw=green!50!black] (triv) at (0,0.0) {$\{0\}$（1 阶）};

\draw[dmSGLatLine] (g) -- (h4);
\draw[dmSGLatLine] (g) -- (h6);
\draw[dmSGLatLine] (h4) -- (h2);
\draw[dmSGLatLine] (h6) -- (h2);
\draw[dmSGLatLine] (h6) -- (h3);
\draw[dmSGLatLine] (h2) -- (triv);
\draw[dmSGLatLine] (h3) -- (triv);

\node at (0, -1) {\fangsong $\mathbb{Z}_{12}$ 的子群格（Hasse 图）：每个因子对应一个子群};
\end{tikzpicture}
\end{center}

\subsection{正规子群与商群}

\begin{mydefinition}{}
子群 $N \leq G$ 称为\textbf{正规子群}，若对任意 $g \in G$：$gNg^{-1} = N$（即 $\forall g\in G,\forall n\in N,\; gng^{-1}\in N$）。记作 $N \trianglelefteq G$。
\end{mydefinition}

\begin{itemize}
    \item Abel 群的任何子群都是正规子群。
    \item 对任意同态 $\varphi: G \to H$，$\ker\varphi \trianglelefteq G$。
    \item 指数为 2 的子群必为正规子群。
\end{itemize}

\begin{mydefinition}{商群}
设 $N \trianglelefteq G$。在左陪集的集合 $G/N = \{aN \mid a \in G\}$ 上定义运算：
\[
(aN)*(bN) = (ab)N
\]
则 $(G/N, *)$ 构成群，称为 $G$ 关于 $N$ 的\textbf{商群}。
若 $G$ 有限，$|G/N| = |G|/|N| = [G:N]$。
\end{mydefinition}

例如 $n\mathbb{Z} \trianglelefteq \mathbb{Z}$，商群 $\mathbb{Z}/n\mathbb{Z} \cong \mathbb{Z}_n$。

\subsection{群同态与同构定理}

\begin{mydefinition}{}
设 $(G,*)$ 和 $(H,\circ)$ 是群。映射 $\varphi: G\to H$ 称为\textbf{群同态}，若：
\[
\forall a,b\in G,\; \varphi(a*b) = \varphi(a)\circ \varphi(b)
\]
\end{mydefinition}

\begin{mydefinition}{核与像}
\begin{itemize}
    \item \textbf{核}：$\ker\varphi = \{g\in G \mid \varphi(g)=e_H\}$，$\ker\varphi \trianglelefteq G$。
    \item \textbf{像}：$\operatorname{im}\varphi = \{\varphi(g) \mid g\in G\}$，$\operatorname{im}\varphi \leq H$。
\end{itemize}
\end{mydefinition}

$\varphi$ 是单射当且仅当 $\ker\varphi = \{e_G\}$。若 $\varphi$ 是双射，则称 $\varphi$ 为\textbf{同构}，记作 $G \cong H$。

\begin{myTheorem}{第一同构定理}
设 $\varphi: G \to H$ 是群同态，则
\[
\boxed{G / \ker\varphi \;\cong\; \operatorname{im}\varphi}
\]
\end{myTheorem}

\begin{myTheorem}{Cayley 定理}
任意群 $G$ 同构于某个对称群的子群。特别地，若 $|G|=n$，则 $G$ 同构于 $S_n$ 的某个子群。
\end{myTheorem}

\section{置换群}

\subsection{置换的表示}

\begin{mydefinition}{}
有限集合 $X = \{1,2,\dots,n\}$ 上的\textbf{置换}是一个双射 $\sigma: X \to X$。
常用双行表示法：
\[
\sigma = \begin{pmatrix}
1 & 2 & 3 & 4 & 5 \\
3 & 1 & 5 & 2 & 4
\end{pmatrix}
\]
\end{mydefinition}

\begin{mydefinition}{轮换}
长度为 $k$ 的\textbf{轮换} $(a_1\;a_2\;\cdots\;a_k)$ 表示置换：
\[
a_1\mapsto a_2,\; a_2\mapsto a_3,\;\dots,\; a_{k-1}\mapsto a_k,\; a_k\mapsto a_1
\]
其余元素不动。$k=2$ 时为\textbf{对换}。
\end{mydefinition}

\begin{myTheorem}{轮换分解定理}
每个置换可以唯一（不计轮换次序和循环移位）地分解为不相交轮换的乘积。
\end{myTheorem}

\begin{example}
$\begin{pmatrix}1&2&3&4&5&6\\ 3&1&5&6&2&4\end{pmatrix} = (1\;3\;5\;2)(4\;6)$
\end{example}

\begin{myTheorem}{}
每个置换可以表示为对换的乘积。表示不唯一，但对换个数的\textbf{奇偶性}唯一。
\end{myTheorem}

\subsection{对称群与交错群}

\begin{mydefinition}{}
$n$ 次\textbf{对称群} $S_n$：$n$ 元集合上所有置换关于复合构成的群。$|S_n| = n!$。
\end{mydefinition}

\begin{itemize}
    \item \textbf{偶置换}：可分解为偶数个对换的乘积。
    \item \textbf{奇置换}：可分解为奇数个对换的乘积。
    \item \textbf{交错群} $A_n$：$S_n$ 中所有偶置换构成的子群。$|A_n| = \frac{n!}{2}$（$n\geq 2$）。$A_n \trianglelefteq S_n$。
\end{itemize}

$S_3 \cong D_3$（6 阶非 Abel 群）;$A_n$（$n \geq 5$）是单群（无非平凡正规子群）。

\section{环与域}

\subsection{环的定义}

\begin{mydefinition}{}
代数系统 $(R,+,\cdot)$ 称为\textbf{环}，若：
\begin{enumerate}
    \item $(R,+)$ 是 Abel 群（加法幺元记作 $0$，$a$ 的加法逆元记作 $-a$）。
    \item $(R,\cdot)$ 是半群（乘法满足封闭性和结合律）。
    \item \textbf{分配律}成立：$a\cdot(b+c) = a\cdot b + a\cdot c$，$(a+b)\cdot c = a\cdot c + b\cdot c$。
\end{enumerate}
\end{mydefinition}

环中的基本计算：$0\cdot a = a\cdot 0 = 0$，$(-a)\cdot b = a\cdot (-b) = -(a\cdot b)$，$(-a)\cdot (-b) = a\cdot b$。

\subsection{环的类型}

\begin{center}
\begin{tabular}{c|c}
\hline
\textbf{环的类型} & \textbf{额外条件} \\
\hline
交换环 & 乘法满足交换律 \\
含幺环 & 乘法有幺元 $1$，且 $1\neq 0$ \\
无零因子环 & 若 $a\cdot b = 0$，则 $a=0$ 或 $b=0$ \\
整环 & 交换含幺环 + 无零因子 \\
除环（体） & 含幺环 + 每个非零元有乘法逆元 \\
域 & 交换的除环 \\
\hline
\end{tabular}
\end{center}

\[
\boxed{\text{域} \;\subset\; \text{整环} \;\subset\; \text{交换含幺环} \;\subset\; \text{环}}
\]
\[
\boxed{\text{域} \;\subset\; \text{除环} \;\subset\; \text{含幺环}}
\]

\begin{mydefinition}{零因子}
环 $R$ 中，非零元 $a$ 称为\textbf{零因子}，若存在 $b \neq 0$ 使得 $a\cdot b = 0$（或 $b\cdot a = 0$）。在 $\mathbb{Z}_6$ 中，$\bar{2}\cdot\bar{3} = \bar{0}$，故 $\bar{2}$ 和 $\bar{3}$ 是零因子。
\end{mydefinition}

\begin{myTheorem}{}
$\mathbb{Z}_n$ 是整环 $\iff$ $\mathbb{Z}_n$ 是域 $\iff$ $n$ 是素数。有限整环必为域。
\end{myTheorem}

常见例子：
\[
\begin{array}{c|c}
\hline
\text{结构} & \text{例子} \\
\hline
\text{非交换环} & M_n(\mathbb{R}),\text{ 四元数 }\mathbb{H} \\
\text{整环（非域）} & \mathbb{Z},\; \mathbb{R}[x]\text{（多项式环）} \\
\text{除环（非域）} & \text{四元数 }\mathbb{H} \\
\text{域} & \mathbb{Q},\mathbb{R},\mathbb{C},\mathbb{Z}_p \\
\hline
\end{array}
\]

\subsection{子环与理想}

\begin{mydefinition}{}
环 $R$ 的子环 $I$ 称为\textbf{理想}，若 $\forall r\in R,\forall i\in I,\; r\cdot i\in I$ 且 $i\cdot r\in I$（吸收性）。
\end{mydefinition}

例如：$n\mathbb{Z} = \{nk \mid k\in\mathbb{Z}\}$ 是 $\mathbb{Z}$ 的理想。多项式环 $\mathbb{R}[x]$ 中以 $0$ 为根的多项式集合也是理想。

\subsection{域}

\begin{mydefinition}{}
代数系统 $(F,+,\cdot)$ 称为\textbf{域}，若：
\begin{enumerate}
    \item $(F,+)$ 是 Abel 群（幺元 $0$）
    \item $(F\setminus\{0\}, \cdot)$ 是 Abel 群（幺元 $1$）
    \item 乘法对加法满足分配律
\end{enumerate}
\end{mydefinition}

\begin{myTheorem}{域的特征}
域 $F$ 的\textbf{特征} $\operatorname{char}(F)$ 是满足 $n\cdot 1=0$ 的最小正整数 $n$;若无此类 $n$，则 $\operatorname{char}(F)=0$。域的特征要么为 0，要么为素数。
\end{myTheorem}

\subsection{有限域}

\begin{myTheorem}{}
有限域的阶（元素个数）必为 $p^n$，其中 $p$ 是素数（域的特征），$n$ 是正整数。
\end{myTheorem}

对每个素数幂 $q = p^n$，存在唯一（同构意义下）的 $q$ 元有限域，记作 $\operatorname{GF}(q)$ 或 $\mathbb{F}_q$。$\operatorname{GF}(p) \cong \mathbb{Z}_p$。$\operatorname{GF}(p^n)$ 的乘法群是 $p^n-1$ 阶\textbf{循环群}。

例如构造 $\operatorname{GF}(4) \cong \mathbb{Z}_2[x] / (x^2+x+1)$，元素为 $\{0, 1, \alpha, \alpha+1\}$，其中 $\alpha^2+\alpha+1=0$（在 $\mathbb{Z}_2$ 中）。

\section{格}

\subsection{格作为偏序集}

\begin{mydefinition}{}
偏序集 $(L, \preceq)$ 称为\textbf{格}，若任意两个元素 $a,b\in L$ 都有：
\begin{itemize}
    \item \textbf{最小上界（并）} $a \vee b$：满足 $a \preceq a\vee b$，$b \preceq a\vee b$，且对任意上界 $c$ 有 $a\vee b \preceq c$。
    \item \textbf{最大下界（交）} $a \wedge b$：满足 $a\wedge b \preceq a$，$a\wedge b \preceq b$，且对任意下界 $c$ 有 $c \preceq a\wedge b$。
\end{itemize}
\end{mydefinition}

\begin{center}
\begin{tikzpicture}[
    dmLatticeNode/.style={circle, draw=blue!60, fill=blue!10, minimum size=7mm, inner sep=0pt},
    dmLatLine/.style={thick, draw=blue!50}
]
% Lattice of subsets {1,2,3} as Hasse diagram
\node[dmLatticeNode] (top) at (0,3.5) {$\{1,2,3\}$};
\node[dmLatticeNode] (s12) at (-1.5,2) {$\{1,2\}$};
\node[dmLatticeNode] (s13) at (0,2) {$\{1,3\}$};
\node[dmLatticeNode] (s23) at (1.5,2) {$\{2,3\}$};
\node[dmLatticeNode] (s1) at (-1.5,0.5) {$\{1\}$};
\node[dmLatticeNode] (s2) at (0,0.5) {$\{2\}$};
\node[dmLatticeNode] (s3) at (1.5,0.5) {$\{3\}$};
\node[dmLatticeNode] (bot) at (0,-1) {$\emptyset$};

\draw[dmLatLine] (top) -- (s12);
\draw[dmLatLine] (top) -- (s13);
\draw[dmLatLine] (top) -- (s23);
\draw[dmLatLine] (s12) -- (s1);
\draw[dmLatLine] (s12) -- (s2);
\draw[dmLatLine] (s13) -- (s1);
\draw[dmLatLine] (s13) -- (s3);
\draw[dmLatLine] (s23) -- (s2);
\draw[dmLatLine] (s23) -- (s3);
\draw[dmLatLine] (s1) -- (bot);
\draw[dmLatLine] (s2) -- (bot);
\draw[dmLatLine] (s3) -- (bot);

\node at (0, -1.8) {\fangsong $(\mathcal{P}(\{1,2,3\}), \subseteq)$ 的 Hasse 图：布尔格 $2^3$};
\end{tikzpicture}
\end{center}

\subsection{格作为代数结构}

\begin{mydefinition}{}
代数系统 $(L, \vee, \wedge)$ 称为\textbf{格}，若对任意 $a,b,c\in L$：
\begin{enumerate}
    \item 交换律：$a\vee b = b\vee a$，$a\wedge b = b\wedge a$
    \item 结合律：$(a\vee b)\vee c = a\vee (b\vee c)$，$(a\wedge b)\wedge c = a\wedge (b\wedge c)$
    \item 吸收律：$a\vee (a\wedge b) = a$，$a\wedge (a\vee b) = a$
\end{enumerate}
\end{mydefinition}

两种定义等价：通过 $a\preceq b \iff a\wedge b=a \iff a\vee b=b$ 建立联系。

\subsection{特殊格}

\begin{center}
\begin{tabular}{c|c}
\hline
\text{类型} & \text{额外条件} \\
\hline
有界格 & 存在最大元 $1$ 和最小元 $0$ \\
分配格 & $a\wedge(b\vee c) = (a\wedge b)\vee(a\wedge c)$ \\
有补格 & 有界 + 每个元素有补元 $a'$（$a\vee a'=1, a\wedge a'=0$）\\
布尔格 & 有补分配格（即布尔代数）\\
\hline
\end{tabular}
\end{center}

\begin{myTheorem}{Birkhoff 刻画}
一个格是分配格当且仅当它不包含同构于 $N_5$（五边形格）或 $M_3$（菱形格）的子格。
\end{myTheorem}

\section{编程实现}

\begin{mycode}{模 $n$ 剩余类群}
\begin{lstlisting}[language=Python]
import math

class Z_n:
    def __init__(self, n):
        self.n = n

    def add(self, a, b):
        return (a + b) % self.n

    def mul(self, a, b):
        return (a * b) % self.n

    def inv_add(self, a):
        return (-a) % self.n

    def inv_mul(self, a):
        def egcd(a, b):
            if b == 0: return (a, 1, 0)
            g, x, y = egcd(b, a % b)
            return (g, y, x - (a // b) * y)
        g, x, _ = egcd(a, self.n)
        if g != 1:
            raise ValueError(f"{a} 与 {self.n} 不互质")
        return x % self.n

    def mul_group(self):
        return {a for a in range(1, self.n)
                if math.gcd(a, self.n) == 1}

    def is_field(self):
        if self.n <= 1: return False
        return all(self.n % i != 0
                   for i in range(2, int(math.isqrt(self.n)) + 1))

# 示例
z7 = Z_n(7)
print(f"Z_7 是域？{z7.is_field()}")
print(f"乘法群元素：{z7.mul_group()}")  # {1,2,3,4,5,6}
print(f"3 的乘法逆元 mod 7：{z7.inv_mul(3)}")  # 5
\end{lstlisting}
\end{mycode}

\begin{mycode}{置换的运算与分解}
\begin{lstlisting}[language=Python]
import math

class Permutation:
    def __init__(self, mapping):
        self.n = max(mapping.keys(), default=0)
        self.mapping = {i: mapping.get(i, i)
                        for i in range(1, self.n + 1)}

    def __call__(self, x):
        return self.mapping.get(x, x)

    def __mul__(self, other):
        new_map = {i: self(other(i))
                   for i in range(1, max(self.n, other.n) + 1)}
        return Permutation(new_map)

    def inverse(self):
        inv_map = {v: k for k, v in self.mapping.items()}
        return Permutation(inv_map)

    def to_cycles(self):
        visited = set()
        cycles = []
        for start in range(1, self.n + 1):
            if start in visited:
                continue
            cycle = []
            cur = start
            while cur not in visited:
                visited.add(cur)
                cycle.append(cur)
                cur = self.mapping[cur]
            if len(cycle) > 1:
                cycles.append(tuple(cycle))
        return cycles

    def parity(self):
        trans = 0
        for c in self.to_cycles():
            trans += len(c) - 1
        return trans % 2  # 0=偶, 1=奇

    def order(self):
        lengths = [len(c) for c in self.to_cycles()]
        if not lengths:
            return 1
        result = lengths[0]
        for L in lengths[1:]:
            result = result * L // math.gcd(result, L)
        return result

    def __repr__(self):
        cycles = self.to_cycles()
        if not cycles:
            return "id"
        return " ".join("(" + " ".join(str(x) for x in c) + ")"
                        for c in cycles)

# 示例
sigma = Permutation({1:2, 2:3, 3:1, 4:5, 5:4})
print(f"sigma = {sigma}")       # (1 2 3) (4 5)
print(f"奇偶性: {sigma.parity()}")  # 1 (奇)
print(f"阶: {sigma.order()}")       # 6
print(f"逆: {sigma.inverse()}")     # (1 3 2) (4 5)
\end{lstlisting}
\end{mycode}

\begin{mycode}{循环群的子群分析}
\begin{lstlisting}[language=Python]
import math

def generators_of_Z_n(n):
    return [k for k in range(1, n) if math.gcd(k, n) == 1]

def order_of_element(n, a):
    if a % n == 0:
        return 1
    return n // math.gcd(a, n)

def subgroups_of_Z_n(n):
    subgroups = {}
    for d in range(1, n + 1):
        if n % d == 0:
            g = n // d
            subgroup = sorted((k * g) % n for k in range(d))
            subgroups[d] = subgroup
    return subgroups

def cyclic_group_info(n):
    print(f"Z_{n} 是 {n} 阶循环群")
    gens = generators_of_Z_n(n)
    print(f"  生成元（φ({n})={len(gens)} 个）：{gens}")
    print("  子群：")
    for order, elements in subgroups_of_Z_n(n).items():
        print(f"    {order} 阶：{elements}")

cyclic_group_info(12)
# Z_12 的子群对应 12 的每个因子
\end{lstlisting}
\end{mycode}

\section{本章小结}

\begin{myTheorem}{代数结构体系总览}
\[
\begin{array}{c|l|l}
\hline
\textbf{结构} & \textbf{定义关键条件} & \textbf{经典例子} \\
\hline
\text{半群} & \text{封闭 + 结合} & (\mathbb{N},+), (\Sigma^+,\circ) \\
\text{幺半群} & \text{半群 + 幺元} & (\mathbb{N}_0,+), (\Sigma^*,\circ) \\
\text{群} & \text{幺半群 + 逆元} & (\mathbb{Z},+), S_n, \operatorname{GL}_n(\mathbb{R}) \\
\text{Abel 群} & \text{群 + 交换律} & (\mathbb{Z},+), (\mathbb{Z}_n,\oplus) \\
\text{环} & \text{Abel 加法群 + 乘法半群 + 分配律} & (\mathbb{Z},+,\times) \\
\text{整环} & \text{交换含幺环 + 无零因子} & \mathbb{Z}, \mathbb{Z}[x] \\
\text{域} & \text{交换除环} & \mathbb{Q},\mathbb{R},\mathbb{C},\mathbb{Z}_p \\
\text{格} & \text{任意两元有上下确界} & (\mathcal{P}(X),\subseteq) \\
\hline
\end{array}
\]
\end{myTheorem}

\textbf{五大核心定理}：

\begin{center}
\begin{tabular}{c|l}
\hline
\textbf{定理} & \textbf{内容} \\
\hline
拉格朗日定理 & 子群阶整除群阶;元素阶整除群阶 \\
Cayley 定理 & 每个群同构于某个对称群的子群 \\
第一同构定理 & $G/\ker\varphi \cong \operatorname{im}\varphi$ \\
循环群结构定理 & $n$ 阶循环群对每个 $d\mid n$ 恰有一个 $d$ 阶子群 \\
有限整环即域 & 有限整环必为域 \\
\hline
\end{tabular}
\end{center}
